\documentclass[aspectratio=169,12pt,t]{beamer}
\usepackage{graphicx}
\setbeameroption{hide notes}
\setbeamertemplate{note page}[plain]
\usepackage{listings}

% header.tex: boring LaTeX/Beamer details + macros

% get rid of junk
\usetheme{default}
\beamertemplatenavigationsymbolsempty
\hypersetup{pdfpagemode=UseNone} % don't show bookmarks on initial view


% font
\usepackage{fontspec}
\setsansfont
  [ ExternalLocation = fonts/ ,
    UprightFont = *-regular ,
    BoldFont = *-bold ,
    ItalicFont = *-italic ,
    BoldItalicFont = *-bolditalic ]{texgyreheros}
\setbeamerfont{note page}{family*=pplx,size=\footnotesize} % Palatino for notes
% "TeX Gyre Heros can be used as a replacement for Helvetica"
% I've placed them in fonts/; alternatively you can install them
% permanently on your system as follows:
%     Download http://www.gust.org.pl/projects/e-foundry/tex-gyre/heros/qhv2.004otf.zip
%     In Unix, unzip it into ~/.fonts
%     In Mac, unzip it, double-click the .otf files, and install using "FontBook"

% named colors
\definecolor{offwhite}{RGB}{255,250,240}
\definecolor{gray}{RGB}{155,155,155}

\ifx\notescolors\undefined % slides (made the same colors as notes)
  \definecolor{background}{RGB}{255,255,255}
  \definecolor{foreground}{RGB}{24,24,24}
  \definecolor{title}{RGB}{27,94,134}
  \definecolor{subtitle}{RGB}{22,175,124}
  \definecolor{hilit}{RGB}{122,0,128}
  \definecolor{vhilit}{RGB}{255,0,128}
  \definecolor{codehilit}{RGB}{24,24,24}
  \definecolor{lolit}{RGB}{95,95,95}
\else % notes
  \definecolor{background}{RGB}{255,255,255}
  \definecolor{foreground}{RGB}{24,24,24}
  \definecolor{title}{RGB}{27,94,134}
  \definecolor{subtitle}{RGB}{22,175,124}
  \definecolor{hilit}{RGB}{122,0,128}
  \definecolor{vhilit}{RGB}{255,0,128}
  \definecolor{codehilit}{RGB}{24,24,24}
  \definecolor{lolit}{RGB}{95,95,95}
\fi
\definecolor{nhilit}{RGB}{128,0,128}  % hilit color in notes
\definecolor{nvhilit}{RGB}{255,0,128} % vhilit for notes

\newcommand{\hilit}{\color{hilit}}
\newcommand{\vhilit}{\color{vhilit}}
\newcommand{\nhilit}{\color{nhilit}}
\newcommand{\nvhilit}{\color{nvhilit}}
\newcommand{\lolit}{\color{lolit}}

% use those colors
\setbeamercolor{titlelike}{fg=title}
\setbeamercolor{subtitle}{fg=subtitle}
\setbeamercolor{institute}{fg=lolit}
\setbeamercolor{normal text}{fg=foreground,bg=background}
\setbeamercolor{item}{fg=foreground} % color of bullets
\setbeamercolor{subitem}{fg=lolit}
\setbeamercolor{itemize/enumerate subbody}{fg=lolit}
\setbeamertemplate{itemize subitem}{{\textendash}}
\setbeamerfont{itemize/enumerate subbody}{size=\footnotesize}
\setbeamerfont{itemize/enumerate subitem}{size=\footnotesize}

% page number
\setbeamertemplate{footline}{%
    \raisebox{5pt}{\makebox[\paperwidth]{\hfill\makebox[20pt]{\lolit
          \scriptsize\insertframenumber}}}\hspace*{5pt}}

% add a bit of space at the top of the notes page
\addtobeamertemplate{note page}{\setlength{\parskip}{12pt}}

% default link color
\hypersetup{colorlinks, urlcolor={hilit}}

\ifx\notescolors\undefined % slides
  % set up listing environment
  \lstset{language=bash,
          basicstyle=\ttfamily\scriptsize,
          frame=single,
          commentstyle=,
          backgroundcolor=\color{offwhite},
          showspaces=false,
          showstringspaces=false
          }
\else % notes
  \lstset{language=bash,
          basicstyle=\ttfamily\scriptsize,
          frame=single,
          commentstyle=,
          backgroundcolor=\color{offwhite},
          showspaces=false,
          showstringspaces=false
          }
\fi

% a few macros
\newcommand{\bi}{\begin{itemize}}
\newcommand{\bbi}{\vspace{24pt} \begin{itemize} \itemsep8pt}
\newcommand{\ei}{\end{itemize}}
\newcommand{\ig}{\includegraphics}
\newcommand{\subt}[1]{{\footnotesize \color{subtitle} {#1}}}
\newcommand{\ttsm}{\tt \small}
\newcommand{\ttfn}{\tt \footnotesize}
\newcommand{\figh}[2]{\centerline{\includegraphics[height=#2\textheight]{#1}}}
\newcommand{\figw}[2]{\centerline{\includegraphics[width=#2\textwidth]{#1}}}


%%%%%%%%%%%%%%%%%%%%%%%%%%%%%%%%%%%%%%%%%%%%%%%%%%%%%%%%%%%%%%%%%%%%%%
% end of header
%%%%%%%%%%%%%%%%%%%%%%%%%%%%%%%%%%%%%%%%%%%%%%%%%%%%%%%%%%%%%%%%%%%%%%

% title info
\title{The culture of open scholarship}
\author{\href{https://kbroman.org}{Karl Broman}}
\institute{Biostatistics \& Medical Informatics \\ Univ.\ Wisconsin{\textendash}Madison}
\date{\href{https://kbroman.org}{\tt \scriptsize \color{foreground} kbroman.org}
\\[2pt]
\scriptsize {\lolit Slides:} \href{https://bit.ly/broman2023}{\tt \scriptsize
  \color{foreground} bit.ly/broman2023}
}


\begin{document}

% title slide
{
\setbeamertemplate{footline}{} % no page number here
\frame{
  \titlepage

  \vfill \hfill \includegraphics[height=6mm]{Figs/cc-zero.png} \vspace*{-1cm}


  \note{
    These are slides for a 10-15 min talk on open scholarship for the Big 10
    Academic Alliance libraries, in Jan 2023.

    Slides (pdf): https://kbroman.org/Talk\_Big10Libs/oa2023.pdf

    Slides with notes (pdf): https://kbroman.org/Talk\_Big10Libs/oa2023\_notes.pdf

    Source: https://github.com/kbroman/Talk\_Big10Libs
  }
} }



\begin{frame}[c]{}

\hfill \begin{minipage}{108mm}

{\Large \color{title}

  open access

\bigskip

  open educational resources

\bigskip

  open source

\bigskip

  open science

}
\end{minipage}

  \note{
      In thinking about Open Scholarship, I'm generally thinking of four
      things: open access publications, open educational resources,
      open source software, and open science (by which I mean open
      data, methods, and materials).

      I'm going to focus mostly on Open Access publications, and on
      academic scientists.

      Concerns in the humanities can be quite different, and I'm going
      to focus on what I know best, which is the situation in science.
  }
\end{frame}


\begin{frame}{About me}

  \vspace{-12pt}

  \bbi
\item Applied statistician working in genetics
\item Co-author on 170 papers and 1 book
\item Reviewer for 90 different journals
\item Formerly
  \bi
\item Associate Editor and Senior Editor at \emph{Genetics}
\item Associate Editor at \emph{Biostatistics}
\item Associate Editor at \emph{Journal of the Americal Statistical
Association}
\item Academic Editor at \emph{PeerJ}
\item Editorial Board member of \emph{BMC Biology}
  \ei
  \ei

  \vspace{30pt}
  \hfill
  \href{https://kbroman.org/broman_cv.pdf}{\tt \footnotesize \lolit kbroman.org/broman{\_}cv.pdf}

  \note{
     I thought I should say a bit about my experience with the
     publication process.

     My research spans genetics and statistics, but the majority of my
     publications are in the genetics or biology literature; I only
     have a couple of papers in statistics journals.

     I've reviewed for a lot of different journals, and I spent a
     decade as an editor for the society journal, Genetics.
  }
\end{frame}



\begin{frame}[c]{}

  \vspace{24pt}

  \centerline{\fbox{\parbox[c]{\textwidth}{
  \figw{Figs/dataorg_paper.png}{1.0}
  }}}

  \vfill
  \hfill
  \href{https://doi.org/gdz6cm}{\tt \footnotesize \lolit doi.org/gdz6cm}

  \note{
    I thought I'd start with a story about the paper of mine.

    This isn't a typical research paper, but really more of a
    tutorial, or really a screed.
  }
\end{frame}


\begin{frame}[c]{}
  \figh{Figs/dataorg_fig5c.png}{0.9}

  \note{
    The paper concerns how to organize data in spreadsheets. And
    really that spreadsheets shouldn't be organized like this example,
    but rather as a rectangle with the columns being the measured
    variables and one row per subject, and a single header row.
  }
\end{frame}


\begin{frame}[c]{}
  \figh{Figs/amstat_mostread.png}{0.9}

  \note{
    It was published in the American Statistician, a society journal,
    and is third-most downloaded paper in that journal, after two
    papers about P-values.
  }
\end{frame}



\begin{frame}[c]{}
  \vspace{8pt}

  \centerline{\fbox{\parbox[c]{0.72\textwidth}{
  \figh{Figs/dataorg_website.png}{0.9}
  }}}

  \vfill
  \hfill
  \href{https://kbroman.org/dataorg}{\tt \footnotesize \lolit kbroman.org/dataorg}

  \note{
    Before it was a paper, it was just a website. It's one of several
    short tutorials that I've written, about things related to
    reproducible research and open science.
  }
\end{frame}


\begin{frame}[c]{}
  \vspace{8pt}

  \centerline{\fbox{\parbox[c]{0.8\textwidth}{
  \figh{Figs/dataorg_preprint.png}{0.9}
  }}}

  \vfill
  \hfill
  \href{https://doi.org/10.7287/peerj.preprints.3183v2}{\tt \footnotesize \lolit doi.org/10.7287/peerj.preprints.3183v2}

  \note{
    In addition, the submitted manuscript is openly available at PeerJ
    Preprints.
  }
\end{frame}


\begin{frame}[c]{}
  \figh{Figs/dataorg_invoice.png}{0.9}

  \note{
    Nevertheless, I paid nearly \$3000 to have the published paper
    available open access.

    I struggled with the decision of whether to pay this fee. But I'm
    glad that I did.

    I think audience for this paper was made much more widespread by
    being a formal paper (rather than a preprint, and before that
    basically a blog post).
  }
\end{frame}


\begin{frame}[c]{}
  \vspace{8pt}

  \centerline{\fbox{\parbox[c]{0.5\textwidth}{
  \figh{Figs/dataorg_funfacts.png}{0.9}
  }}}

  \vfill
  \hfill
  \href{https://bit.ly/3sIRtVY}{\tt \footnotesize \lolit bit.ly/3sIRtVY}

  \note{
    If you want to more about this paper, see my twitter thread on 10
    fun facts about the paper.
  }
\end{frame}


\begin{frame}[c]{Personal timeline}

\figw{Figs/timeline.pdf}{1.0}

  \note{
  }
\end{frame}


\begin{frame}{What has changed?}

  \bbi
\item Rise of preprints in biology
\item Rise of \emph{Nature Communications}
\item PubMed Central: expansion, with no embargo
\item Emphasis on computational reproducibility
  \ei

  \note{
  }
\end{frame}


\begin{frame}{What hasn't changed?}

  \bbi
\item Attachment to Impact Factor
\item Attachment to Science, Nature, and Cell
\item Researchers don't read much
  \ei

  \note{
    The journal Genetics recently sent out a newsletter that
    included an announcement of new associate editors. The bio for one
    of them had ``...including more than 30 articles in Nature, Science,
    Cell, and Nature Genetics.''
  }
\end{frame}


\begin{frame}{Barriers to OA}

  \bbi
\item Most scientists don't seem to care
\item Cost
  \ei

  \note{
  }
\end{frame}



\end{document}
